%%% FIELD worked-trajectory-construction
For a deterministic path \(w\in\mathcal A^8\), intervene to set \(W=w\), evaluate the eight potential outcomes along that path, and integrate only the latent variables drawn from \(\mathsf Q\).  Denote the resulting Lebesgue density on \([-B,B]^8\) by \(\ell_{\mathsf Q}(y;w)\); this is a structural path likelihood and is not the conditional density of \(Y^{\mathrm{obs}}\) given \(W=w\) under the adaptive experiment.  For \(\mathsf Q\in\mathcal Q_8\), define the joint trajectory density with respect to counting measure on \(\mathcal A^8\) times Lebesgue measure by
\[
 P_{\mathsf Q,\pi^\dagger}(w,y)
 =g_{\pi^\dagger}(w,y)\ell_{\mathsf Q}(y;w).
\]

%%% FIELD worked-cancellation-statement
At \(T=8\), every prior \(\mathsf Q\in\mathcal Q_8\) is fixed before assignment and satisfies \(\mathsf Q(\mathcal P_{8,\beta}(B,C))=1\).  For every \(w\in\mathcal A^8\) and Lebesgue-almost every \(y\in[-B,B]^8\), the joint density under the outcome-adaptive policy \(\pi^\dagger\) factors as
\[
 P_{\mathsf Q,\pi^\dagger}(w,y)
 =\left[\frac12\prod_{t=2}^8p_t^\dagger(y)^{w_t}
 \{1-p_t^\dagger(y)\}^{1-w_t}\right]\ell_{\mathsf Q}(y;w).
\]
Moreover, \(P_{\mathsf Q,\pi^\dagger}\ll P_{Q_0^{(8)},\pi^\dagger}\), and, \(P_{Q_0^{(8)},\pi^\dagger}\)-almost surely,
\[
 \frac{dP_{\mathsf Q,\pi^\dagger}}
 {dP_{Q_0^{(8)},\pi^\dagger}}(w,y)
 =\frac{\ell_{\mathsf Q}(y;w)}{\ell_{Q_0^{(8)}}(y;w)}.
\]
Thus the adaptive assignment factor cancels exactly, without identifying \(\ell_{\mathsf Q}(y;w)\) with an observational conditional density.  Along either alternating path, \(h_{t,a}^{(8)}(w)=0\) for every realized arm and every \(t\ge2\), whereas each arm-specific band prior moves the prior mean of \(\theta_8\) above the null mean by exactly \(7\eta_8c_1/12\).  Along either constant path, the four opposite-arm signs under \(Q_{8,0}^{\mathrm{blk}}\) never enter the likelihood, and their contribution alone gives posterior Bayes variance at least \(\delta_8^2/9\).

%%% FIELD worked-cancellation-proof
Put \(\phi(u)=1-4u^2/B^2\), so that \(\psi_q(u)=u+q\phi(u)\).  For \(u\in[-B/2,B/2]\) and \(\lvert q\rvert\le B/4\),
\[
 0\le\phi(u)\le1,
 \qquad
 \psi_q'(u)=1-\frac{8qu}{B^2}\ge0,
 \qquad
 \psi_q(-B/2)=-B/2,
 \qquad
 \psi_q(B/2)=B/2.
\]
The derivative can vanish only at an endpoint when \(\lvert q\rvert=B/4\); hence \(\psi_q\) is increasing and maps \([-B/2,B/2]\) onto itself.  In the block priors, \(\lvert q\rvert=\delta_8=B/8\).  In the band priors,
\[
 0\le h_{t,a}^{(8)}\le c_1,
 \qquad
 \eta_8c_1\le \frac{Bc_1}{4(c_1+B)}<\frac B4.
\]
Thus every constructed outcome lies in \([-B/2,B/2]\subset[-B,B]\), establishing the support condition in \(\mathrm{ass{:}bounded\mbox{-}potential\mbox{-}outcomes}\).

All latent variables in \(\mathrm{def{:}worked\mbox{-}common\mbox{-}null\mbox{-}family}\) are drawn before assignment.  Under \(Q_0^{(8)}\) and either block prior, a potential outcome depends on a history only through its current arm, so its historical difference is zero.  Under a band prior, two histories with the same current arm have band scores satisfying
\[
 \left|h_{t,a}^{(8)}(z)-h_{t,a}^{(8)}(z')\right|
 \le c_1\mathbf 1\{z_{t-1}\ne z'_{t-1}\}.
\]
Indeed, if the common terminal arm is not \(a\), both scores vanish; if it is \(a\), the score is \(c_1\) or zero according as the lag-one coordinate equals \(a\) or not.  Since
\[
 |\psi_q(u)-\psi_{q'}(u)|=|q-q'|\phi(u)\le |q-q'|
\]
and \(\eta_8\le1/2\),
\[
 |Y_t(z)-Y_t(z')|
 \le \eta_8c_1\mathbf 1\{z_{t-1}\ne z'_{t-1}\}
 \le C2^{-(1+\beta)}\mathbf 1\{z_{t-1}\ne z'_{t-1}\}.
\]
This is at most the right side of \(\mathrm{ass{:}polynomial\mbox{-}historical\mbox{-}envelope}\).  Therefore every \(\mathsf Q\in\mathcal Q_8\) is fixed before assignment and is supported on \(\mathcal P_{8,\beta}(B,C)\).

We next prove the likelihood statement from the generative law.  Fix \(w\).  Under the intervention \(W=w\), integrating the prior latents gives the structural density \(\ell_{\mathsf Q}(y;w)\) from \(\mathrm{def{:}worked\mbox{-}trajectory\mbox{-}laws}\).  For any latent realization that produces the pair \((w,y)\), sequential assignment under \(\pi^\dagger\) contributes the factor
\[
 g_{\pi^\dagger}(w,y)
 =\frac12\prod_{t=2}^8p_t^\dagger(y)^{w_t}
 \{1-p_t^\dagger(y)\}^{1-w_t}.
\]
This factor is measurable in \((w,y)\) and therefore factors out of the integration over the latent variables.  Consequently,
\[
 P_{\mathsf Q,\pi^\dagger}(w,y)
 =g_{\pi^\dagger}(w,y)\ell_{\mathsf Q}(y;w).
\]
This calculation does not condition the adaptive experiment on \(W=w\); it first fixes the path structurally and then multiplies by the sequential assignment probability.

For the null, the eight selected coordinates \(U_{t,w_t}\) are independent uniform variables, so
\[
 \ell_{Q_0^{(8)}}(y;w)=B^{-8}\mathbf 1\{y\in[-B/2,B/2]^8\}.
\]
Conditional on every sign in a block prior, or directly under a band prior, the selected coordinates are independent monotone transports of uniform variables.  Their densities are supported on the same cube and are absolutely continuous with respect to Lebesgue measure; mixing over the block signs preserves this property.  It follows that \(\ell_{\mathsf Q}(\,cdot\,;w)\ll\ell_{Q_0^{(8)}}(\,cdot\,;w)\) for every \(w\).  The assignment factors satisfy \(g_{\pi^\dagger}(w,y)>0\) because each factor is \(1/4\) or \(3/4\).  Division of the two joint densities therefore yields
\[
 \frac{dP_{\mathsf Q,\pi^\dagger}}
 {dP_{Q_0^{(8)},\pi^\dagger}}(w,y)
 =\frac{\ell_{\mathsf Q}(y;w)}
 {\ell_{Q_0^{(8)}}(y;w)}
\]
on the null-almost-sure cube, proving exact cancellation.

On an alternating path, if \(w_t=a\), then \(w_{t-1}\ne a\), so \([c_1-2c_1]_+=0\); if \(w_t\ne a\), the leading indicator in \(h_{t,a}^{(8)}\) is zero.  Thus every realized band score is zero.  On the sustained history at arm \(a\), the score equals zero at \(t=1\) and \(c_1\) for \(2\le t\le8\).  If \(U\sim\operatorname{Unif}[-B/2,B/2]\), then
\[
 \mathbb E U=0,
 \qquad
 \mathbb E U^2=B^2/12,
 \qquad
 \mathbb E\phi(U)=1-4\mathbb EU^2/B^2=2/3.
\]
Therefore the perturbed sustained-arm endpoint has mean \((2/3)s_a\eta_8c_1\) in periods \(2,\ldots,8\).  When \(a=1\), this enters the GATE with a plus sign; when \(a=0\), both \(s_0=-1\) and the minus sign in the GATE reverse it.  In either case,
\[
 \mathbb E_{Q_{8,a}^{\mathrm{band}}}\theta_8
 -\mathbb E_{Q_0^{(8)}}\theta_8
 =\frac18\,7\,\frac23\eta_8c_1
 =\frac{7\eta_8c_1}{12}.
\]

Finally, condition under \(Q_{8,0}^{\mathrm{blk}}\) on a constant path \(w_t=a\), its observed outcomes, the offset, and all latent variables other than the four signs \(\xi_{I,1-a}^{(0)}\).  None of these opposite-arm signs occurs in an observed outcome or in a later assignment probability, because the path never uses arm \(1-a\).  They consequently retain their mutually independent Rademacher laws.  The contribution of block \(I\in\mathcal I_{8,0}\) to the conditional variance of \(\theta_8\), up to an irrelevant sign, is
\[
 \frac{\delta_8^2}{64}
 \left\{\sum_{t\in I}\phi(U_{t,1-a})\right\}^2.
\]
The conditional-variance decomposition and independence across blocks give
\[
 \begin{aligned}
 \operatorname{Var}(\theta_8\mid W=w,Y^{\mathrm{obs}})
 &\ge \frac{\delta_8^2}{64}
 \sum_{I\in\mathcal I_{8,0}}
 \mathbb E\left[\left\{\sum_{t\in I}\phi(U_{t,1-a})\right\}^2\right]\\
 &\ge \frac{\delta_8^2}{64}\,4
 \left\{2\mathbb E\phi(U)\right\}^2
 =\frac{\delta_8^2}{9},
 \end{aligned}
\]
where the second inequality is \(\mathbb E[V^2]\ge(\mathbb EV)^2\).  This proves the posterior Bayes-variance assertion.

%%% FIELD multiscale-family-construction
Let \(c_\ell=C(1+\ell)^{-(1+\beta)}\), let
\[
 \mathcal S_L=\{2^j:j\in\mathbb Z_{\ge0},\ 2^j\le L\}\cup\{L\},
 \qquad
 A_L=2^{-(\beta+4)}\{B\wedge(CL^{-\beta})\},
\]
and, for \(b\in\mathcal S_L\), put
\[
 S_b=\sum_{\ell=b}^{2b-1}c_\ell.
\]
Draw the offsets for the dyadic values of \(b\) jointly, before assignment, so that each \(R_b\) is uniform on \(\{0,\ldots,b-1\}\) and \(R_b=R_{2b}\bmod b\) whenever both \(b\) and \(2b\) are dyadic members of \(\mathcal S_L\).  If \(L\) is not dyadic, draw its terminal offset \(R_L\) independently and uniformly.  For each \(b\), form the nonempty intervals
\[
 I_{b,R_b,k}=\{1,\ldots,T\}\cap
 \{kb+R_b+1,\ldots,(k+1)b+R_b\},
 \qquad k\in\mathbb Z.
\]
For a common-null prior \(Q_0^{T,L}\), draw \(U_{t,a}\) independently from \(\operatorname{Unif}[-B/2,B/2]\) for \(1\le t\le T\) and \(a\in\mathcal A\), before assignment, and set \(Y_t(z_{1:t})=U_{t,z_t}\).  For each \(b\in\mathcal S_L\), the randomized-offset block prior \(\overline Q_b^{\mathrm{blk}}\) additionally draws independent Rademacher signs \(\xi_{I,a}\), one for every interval \(I=I_{b,R_b,k}\) and arm \(a\), and sets
\[
 Y_t(z_{1:t})
 =\psi_{(B/16)\xi_{I,z_t}}(U_{t,z_t}),
 \qquad t\in I.
\]
For \(C>0\), every \(a\in\mathcal A\), and every \(b\in\mathcal S_L\), define
\[
 h_{t,a,b}(z_{1:t})
 =\mathbf 1\{z_t=a\}A_L
 \left[
 1-\frac{2}{S_b}
 \sum_{\ell=b}^{\min\{2b-1,t-1\}}c_\ell
 \mathbf 1\{z_{t-\ell}\ne a\}
 \right]_+,
\]
and let the arm-specific band prior \(Q_{a,b}^{\mathrm{band}}\) set
\[
 Y_t(z_{1:t})=\psi_{s_a h_{t,a,b}(z_{1:t})}(U_{t,z_t}).
\]
When \(C=0\), define every band prior to equal \(Q_0^{T,L}\), avoiding division by \(S_b=0\).  The finite family is
\[
 \mathfrak Q_{T,L}
 =\{Q_0^{T,L}\}
 \cup\{\overline Q_b^{\mathrm{blk}}:b\in\mathcal S_L\}
 \cup\{Q_{a,b}^{\mathrm{band}}:a\in\mathcal A,\ b\in\mathcal S_L\}.
\]
The coupled partitions for the dyadic members of \(\mathcal S_L\) are the levels of a randomized dyadic interval tree; the explicit random root offset prevents a path from exploiting a fixed collection of boundaries.  All offsets, signs, and uniforms are drawn before assignment.  Finally define the two normalized path diagnostics
\[
 \overline{\mathsf V}_b(w,r)
 =\frac1{T^2}\sum_{I\in\{I_{b,r,k}\}:\,|\{w_t:t\in I\}|=1}|I|^2
\]
and, when \(C>0\),
\[
 \mathsf I_{a,b}(w)
 =A_L^{-2}\sum_{t=1}^T h_{t,a,b}(w_{1:t})^2.
\]
Set \(\mathsf I_{a,b}=0\) when \(A_L=0\).

%%% FIELD multiscale-payoffs-statement
Fix \(B>0\), \(C\ge0\), \(\beta>0\), \(T\ge4\), and \(1\le L\le T/4\).  Every member of the enumerated family \(\mathfrak Q_{T,L}\) is drawn before assignment and is supported on \(\mathcal P_{T,\beta}(B,C)\).  For every \(\pi\in\mathcal D_T\), every \(b\in\mathcal S_L\), and every \(a\in\mathcal A\), the following exact construction-side bounds hold.

First, the randomized-offset block prior satisfies
\[
 \inf_{\widehat\theta}
 \mathbb E_{\overline Q_b^{\mathrm{blk}},\pi}
 \left[(\widehat\theta-\theta_T(Y))^2\right]
 \ge \frac{B^2}{576}
 \mathbb E_{\overline Q_b^{\mathrm{blk}},\pi}
 [\overline{\mathsf V}_b(W,R_b)].
\]
The common null itself satisfies the policy-uniform floor
\[
 \inf_{\widehat\theta}
 \mathbb E_{Q_0^{T,L},\pi}[(\widehat\theta-\theta_T(Y))^2]
 \ge\frac{B^2}{12T}.
\]
Second, for \(C>0\), the band prior has the exact prior-mean displacement
\[
 \mathbb E_{Q_{a,b}^{\mathrm{band}}}[\theta_T(Y)]
 -\mathbb E_{Q_0^{T,L}}[\theta_T(Y)]
 =\frac{2A_L}{3},
\]
and its reverse structural divergence under the adaptive experiment is bounded by
\[
 D_{\mathrm{KL}}\!\left(
 P_{Q_0^{T,L},\pi}\,\middle\|\,
 P_{Q_{a,b}^{\mathrm{band}},\pi}
 \right)
 \le \frac{16A_L^2}{3B^2}
 \mathbb E_{Q_0^{T,L},\pi}[\mathsf I_{a,b}(W)].
\]
Consequently, its explicit two-prior constrained-risk payoff is
\[
 \begin{aligned}
 &\inf_{\widehat\theta}\max_{Q\in\{Q_0^{T,L},Q_{a,b}^{\mathrm{band}}\}}
 \mathbb E_{Q,\pi}[(\widehat\theta-\theta_T(Y))^2]\\
 &\qquad\ge \frac{A_L^2}{72}
 \left[
 1-\left\{
 \frac{8A_L^2}{3B^2}
 \mathbb E_{Q_0^{T,L},\pi}[\mathsf I_{a,b}(W)]
 \right\}^{1/2}
 -\frac{24B^2}{TA_L^2}
 \right]_+.
 \end{aligned}
\]
For every \(Q\in\mathfrak Q_{T,L}\), the joint density factors as
\[
 P_{Q,\pi}(w,y)=g_\pi(w,y)\ell_Q(y;w),
\]
and wherever the null density is positive,
\[
 \frac{dP_{Q,\pi}}{dP_{Q_0^{T,L},\pi}}(w,y)
 =\frac{\ell_Q(y;w)}{\ell_{Q_0^{T,L}}(y;w)}.
\]
These bounds expose the normalized block posterior-variance payoff and the normalized band information payoff.  They do not assert the missing dyadic stopping/Carleson inequality that would compare the payoffs uniformly over all paths or path laws.

%%% FIELD multiscale-payoffs-proof
All variables defining a prior in \(\mathfrak Q_{T,L}\) are sampled before assignment by construction.  Put \(\phi(u)=1-4u^2/B^2\).  As in the direct calculation for \(\mathrm{lem{:}worked\mbox{-}common\mbox{-}null\mbox{-}cancellation}\), \(\psi_q(u)=u+q\phi(u)\) maps \([-B/2,B/2]\) onto itself whenever \(|q|\le B/4\).  The block amplitude is \(B/16\), and the band amplitude obeys
\[
 0\le h_{t,a,b}\le A_L
 \le 2^{-(\beta+4)}B<B/16.
\]
Thus bounded support holds for every member.

Block-prior outcomes depend on a history only through the terminal arm, so their historical difference is zero.  For a band prior, note first that
\[
 S_b
 =C\sum_{\ell=b}^{2b-1}(1+\ell)^{-(1+\beta)}
 \ge Cb(2b)^{-(1+\beta)}
 =C2^{-(1+\beta)}b^{-\beta}.
\]
Because \(b\le L\),
\[
 \frac{2A_L}{S_b}
 \le
 \frac{2\cdot2^{-(\beta+4)}CL^{-\beta}}
 {C2^{-(1+\beta)}b^{-\beta}}
 \le\frac14.
\]
The positive-part map is one-Lipschitz.  Hence, for histories \(z,z'\) with the same current arm,
\[
 \begin{aligned}
 |h_{t,a,b}(z)-h_{t,a,b}(z')|
 &\le \frac{2A_L}{S_b}
 \sum_{\ell=b}^{\min\{2b-1,t-1\}}c_\ell
 \mathbf 1\{z_{t-\ell}\ne z'_{t-\ell}\}\\
 &\le \sum_{\ell=1}^{t-1}c_\ell
 \mathbf 1\{z_{t-\ell}\ne z'_{t-\ell}\}.
 \end{aligned}
\]
The first line is zero if the common current arm differs from \(a\).  Since \(|\psi_q(u)-\psi_{q'}(u)|\le|q-q'|\), the same upper bound holds for the corresponding potential-outcome difference.  This proves \(\mathrm{ass{:}polynomial\mbox{-}historical\mbox{-}envelope}\), and hence support on \(\mathcal P_{T,\beta}(B,C)\).  When \(C=0\), the band priors equal the null and the conclusion is immediate.

Under \(Q_0^{T,L}\), condition on the complete observed data.  At each time, the unobserved coordinate \(U_{t,1-W_t}\) occurs in neither an observed outcome nor any later adaptive assignment probability, and therefore retains its independent \(\operatorname{Unif}[-B/2,B/2]\) posterior law.  Its coefficient in the GATE is \(1/T\) up to sign.  Since a uniform variable on this interval has variance \(B^2/12\), these \(T\) independent missing coordinates contribute exactly \(B^2/(12T)\) to the conditional posterior variance.  Conditional squared error is posterior variance plus the squared distance from the posterior mean, proving the common-null floor.

We prove the block payoff.  Condition on the offset \(R_b=r\), the complete observed data, and all latent variables except the opposite-arm signs in blocks that are monochromatic along the realized path.  If \(I\) is such a block and its realized arm is \(a_I\), then \(\xi_{I,1-a_I}\) occurs in no observed outcome.  Since later assignment probabilities are functions only of observed past assignments and outcomes, it occurs in no assignment-kernel factor either.  It therefore retains its independent Rademacher posterior law.  Its contribution to the target, up to a sign that does not affect variance, is
\[
 \frac{B}{16T}\xi_{I,1-a_I}
 \sum_{t\in I}\phi(U_{t,1-a_I}).
\]
The conditional-variance identity, followed by \(\mathbb E[V^2]\ge(\mathbb EV)^2\) and \(\mathbb E\phi(U)=2/3\), gives
\[
 \begin{aligned}
 \mathbb E\!\left[
 \operatorname{Var}(\theta_T(Y)\mid W,Y^{\mathrm{obs}},R_b)
 \right]
 &\ge \frac{B^2}{256T^2}
 \mathbb E\!\left[
 \sum_{I:\,|\{W_t:t\in I\}|=1}
 \left\{\sum_{t\in I}\phi(U_{t,1-a_I})\right\}^2
 \right]\\
 &\ge \frac{B^2}{576}
 \mathbb E\!\left[
 \frac1{T^2}\sum_{I:\,|\{W_t:t\in I\}|=1}|I|^2
 \right]\\
 &=\frac{B^2}{576}
 \mathbb E[\overline{\mathsf V}_b(W,R_b)].
 \end{aligned}
\]
Equivalently, the second line follows directly from
\(B^2(2/3)^2/(256T^2)=B^2/(576T^2)\); the normalization by \(T^2\) is exactly the one in \(\overline{\mathsf V}_b\).  For every observed-data estimator, conditional squared error equals posterior variance plus the squared distance from the posterior mean.  Taking the infimum over estimators therefore preserves the lower bound.

For the band displacement, a sustained history at arm \(a\) has no mismatches in the lag band, so \(h_{t,a,b}=A_L\) for every \(t\); the other sustained endpoint is unperturbed.  Since \(\mathbb E\psi_q(U)=q\mathbb E\phi(U)=2q/3\), the sign \(s_1=1\) agrees with the plus sign of the treated endpoint and \(s_0=-1\) cancels the minus sign of the control endpoint.  Every one of the \(T\) summands in \(\theta_T\) is therefore displaced by \(2A_L/3\), proving the stated exact mean shift.

It remains to prove the likelihood and divergence statements.  Given a deterministic path \(w\), integrate the prior latents after structurally setting \(W=w\); call the resulting density \(\ell_Q(y;w)\).  The sequential probability of producing \(w\) when the observed vector is \(y\) is
\[
 g_\pi(w,y)=\prod_{t=1}^T
 \pi_t(w_t\mid w_{1:t-1},y_{1:t-1}).
\]
This factor depends only on \((w,y)\), so it factors out of the latent integration and gives
\[
 P_{Q,\pi}(w,y)=g_\pi(w,y)\ell_Q(y;w).
\]
The monotone transports preserve the null cube and have positive derivative almost everywhere, so every family member is absolutely continuous with respect to the null.  Division gives the Radon--Nikodym identity on the null support, including when some assignment probabilities are zero: paths on which the common factor is zero have zero mass under both laws.

For the quantitative band bound, let \(f_0=1/B\) on \([-B/2,B/2]\), and let \(f_q\) be the density of \(\psi_q(U)\).  With \(d_q(u)=1-8qu/B^2\), change of variables gives
\[
 \frac{f_0(\psi_q(u))}{f_q(\psi_q(u))}=d_q(u).
\]
Using \(\log x\le x-1\), followed by the definition of chi-square divergence,
\[
 \begin{aligned}
 D_{\mathrm{KL}}(f_0\|f_q)
 &\le \int\left(\frac{f_0}{f_q}-1\right)f_0
 =\int\left(\frac{f_0}{f_q}-1\right)^2f_q\\
 &=\mathbb E[(d_q(U)-1)^2]
 =\frac{16q^2}{3B^2},
 \end{aligned}
\]
where \(\mathbb EU^2=B^2/12\).  Under the null, after \(W_t\) is assigned, \(Y_t^{\mathrm{obs}}=U_{t,W_t}\) is uniform and independent of the past.  Under the band prior, its conditional density is \(f_{s_ah_{t,a,b}(W_{1:t})}\).  Expanding the log joint-likelihood ratio into its periodwise factors, the identical assignment kernels cancel at each period.  Iterated expectation and the one-period inequality yield
\[
 \begin{aligned}
 D_{\mathrm{KL}}\!\left(
 P_{Q_0^{T,L},\pi}\|P_{Q_{a,b}^{\mathrm{band}},\pi}
 \right)
 &\le \frac{16}{3B^2}
 \mathbb E_{Q_0^{T,L},\pi}
 \left[\sum_{t=1}^T h_{t,a,b}(W_{1:t})^2\right]\\
 &=\frac{16A_L^2}{3B^2}
 \mathbb E_{Q_0^{T,L},\pi}[\mathsf I_{a,b}(W)].
 \end{aligned}
\]
To convert this divergence into the stated constrained-risk payoff without assuming a fixed target under either prior, write
\[
 \mu_0=\mathbb E_{Q_0^{T,L}}\theta_T(Y)=0,
 \qquad
 \mu_1=\mathbb E_{Q_{a,b}^{\mathrm{band}}}\theta_T(Y)=\frac{2A_L}{3},
 \qquad
 \Delta=\mu_1-\mu_0.
\]
Under the null, independence over \(t\) gives
\[
 \operatorname{Var}_{Q_0^{T,L}}(\theta_T)=\frac{B^2}{6T}.
\]
Under the band prior, the periodwise endpoint contrasts remain independent over \(t\) and lie in \([-B,B]\); each is a difference of two independent variables in \([-B/2,B/2]\), so its variance is at most \(B^2/2\), and
\[
 \operatorname{Var}_{Q_{a,b}^{\mathrm{band}}}(\theta_T)\le\frac{B^2}{2T}.
\]
Let \(E_i=\{|\theta_T-\mu_i|\le\Delta/4\}\) under prior \(i\).  Chebyshev's inequality, obtained directly by applying Markov's inequality to \((\theta_T-\mu_i)^2\), gives
\[
 Q_0^{T,L}(E_0^c)+Q_{a,b}^{\mathrm{band}}(E_1^c)
 \le \frac{24B^2}{TA_L^2}.
\]
For any estimator, test the two priors by choosing the mean closer to \(\widehat\theta\).  On \(E_i\), a testing error forces \(|\widehat\theta-\theta_T|\ge\Delta/4\).  If \(r_i\) denotes its integrated squared risk under prior \(i\), then
\[
 r_0+r_1
 \ge\frac{\Delta^2}{16}
 \left[1-\|P_{Q_0^{T,L},\pi}-P_{Q_{a,b}^{\mathrm{band}},\pi}\|_{\mathrm{TV}}
 -\frac{24B^2}{TA_L^2}\right]_+.
\]
For completeness, the testing inequality in the bracket follows by integrating the smaller of the two observed-data densities.  We also derive the needed total-variation bound.  For probability laws \(P\ll Q\), let \(A=\{dP/dQ\ge1\}\), \(p=P(A)\), \(q=Q(A)\), and \(v=p-q=\|P-Q\|_{\mathrm{TV}}\).  Jensen's inequality under the conditional laws on \(A\) and \(A^c\) gives the two-cell log-sum bound
\[
 D_{\mathrm{KL}}(P\|Q)
 \ge p\log\frac pq+(1-p)\log\frac{1-p}{1-q}.
\]
For fixed \(q\in(0,1)\), the right side, as a function of \(p\), has value and first derivative zero at \(p=q\), while its second derivative is
\(1/\{p(1-p)\}\ge4\).  Twice integrating this lower bound, with the same argument when \(p<q\), proves
\[
 D_{\mathrm{KL}}(P\|Q)\ge2(p-q)^2=2\|P-Q\|_{\mathrm{TV}}^2.
\]
The cases \(q\in\{0,1\}\) follow by a limit or are trivial under \(P\ll Q\).  Applied to the two trajectory laws, this yields
\[
 \|P_{Q_0^{T,L},\pi}-P_{Q_{a,b}^{\mathrm{band}},\pi}\|_{\mathrm{TV}}
 \le\sqrt{D_{\mathrm{KL}}(P_{Q_0^{T,L},\pi}\|P_{Q_{a,b}^{\mathrm{band}},\pi})/2}.
\]
Substituting the divergence bound, using \(\max(r_0,r_1)\ge(r_0+r_1)/2\), and noting \(\Delta^2/32=A_L^2/72\), proves the constrained-risk display.  This proves every asserted construction-side payoff without invoking the unproved dyadic stopping inequality.

%%% FIELD oeq-current-statement
Determine whether there are explicit constants \(\kappa_\beta>0\) and \(T_\beta<\infty\) such that, for every \(T\ge T_\beta\), every \(1\le L\le T/4\), and every \(\pi\in\mathcal D_T\), one can select before assignment a prior \(Q_{\pi,L}\in\mathfrak Q_{T,L}\) satisfying \(Q_{\pi,L}(\mathcal P_{T,\beta}(B,C))=1\) and
\[
 \inf_{\widehat\theta}\mathbb E_{Q_{\pi,L},\pi}
 [ (\widehat\theta-\theta_T(Y))^2]
 \ge\frac{\kappa_\beta}{\log(\mathrm eL)}
 \min\left\{\frac{B^2L}{T},\ B^2\wedge C^2L^{-2\beta}\right\}.
\]
The certificate must be obtained on a common null so that the sequential assignment-kernel factors cancel under outcome adaptation.  Determine separately whether the factor \(1/\log(\mathrm eL)\) can be removed.

%%% FIELD multiscale-open-what
The requested node \(\mathrm{lem{:}dyadic\mbox{-}stopping\mbox{-}occupancy\mbox{-}information}\) remains unproved.  What is needed is a policy-uniform rule, measurable with respect to the known kernel before assignment, selecting either a randomized-offset block prior or an arm-specific band prior so that its Bayes risk is at least
\[
 \frac{\kappa_\beta}{\log(\mathrm eL)}
 \min\left\{\frac{B^2L}{T},\ B^2\wedge C^2L^{-2\beta}\right\}.
\]
Equivalently, a valid charging result must combine the normalized posterior payoff \(\overline{\mathsf V}_b(W,R_b)\) and the normalized structural-information payoff \(\mathsf I_{a,b}(W)\), with all expectations taken under laws for which the comparison is justified.

%%% FIELD multiscale-open-obstruction
The enumerated family and both exact payoffs are now available, but they live under different joint laws.  The block calculation gives \(B^2\mathbb E_{\overline Q_b^{\mathrm{blk}},\pi}[\overline{\mathsf V}_b(W,R_b)]/576\), whereas the band constrained-risk payoff is a decreasing explicit function of \(A_L^2\mathbb E_{Q_0^{T,L},\pi}[\mathsf I_{a,b}(W)]/B^2\).  A deterministic Carleson inequality for a fixed pair \((w,r)\) would compare \(\overline{\mathsf V}_b(w,r)\) and \(\mathsf I_{a,b}(w)\), but it cannot by itself replace the block-law expectation by the null-law expectation.  In particular, although \(R_b\) is uniform before assignment, it need not remain independent of \(W\) under the block prior because its transported observed outcomes can affect later assignments.  The common assignment factor cancels in Radon--Nikodym derivatives, yet the randomized block perturbations accumulate structural information over the horizon, so no uniform bounded-divergence transfer from \(Q_0^{T,L}\) to \(\overline Q_b^{\mathrm{blk}}\) follows.  Thus the missing result is simultaneously a first-separating-node charging inequality and a justified common-measure transfer; assuming either part would assume the crux.

%%% FIELD multiscale-open-attempted
The horizon-eight construction was first repaired by replacing the false observational conditional density with the structural likelihood \(\ell_Q(y;w)\).  The same factorization was then proved for the general finite family.  For every randomized-offset block level, conditioning on the offset and the data gives the exact posterior lower payoff \(B^2\overline{\mathsf V}_b/576\).  For every arm and dyadic band, the endpoint displacement is exactly \(2A_L/3\), and a direct change-of-variables calculation gives the reverse divergence bound \(16A_L^2\mathsf I_{a,b}/(3B^2)\) under the common null.  Decomposing paths into runs and charging a run to its first separating dyadic node handles constant paths, strict alternation, one-switch paths, and homogeneous periodic runs.  The finite-path kill test already rejects the naive dichotomy
\[
 \max_{b\in\{1,2,4\}}\frac{T}{L}
 \mathbb E_r[\overline{\mathsf V}_b(w,r)]\ge\frac12
 \quad\hbox{or}\quad
 \min_{a\in\mathcal A,\,b\in\{1,2,4\}}\mathsf I_{a,b}(w)\le1.
\]
Indeed, at \(\beta=1\), \(T=16\), \(L=4\), and \(w=0000111001110101\), direct enumeration of all offsets gives the first maximum as \(17/64\), while direct evaluation of the lag-band scores gives the second minimum as \(6336489606/5365416001>1\).  Thus a proof needs a genuinely summed, constant-sensitive charging inequality rather than this single-threshold alternative.  Averaging any such pathwise inequality under the null would still leave the block payoff under the wrong law, while averaging under a block prior would leave the band divergence under the wrong law.  No valid inequality found in this attempt closes that change of measure.

%%% FIELD multiscale-open-partial
For the corrected enumerated family, every member is supported on \(\mathcal P_{T,\beta}(B,C)\) and is sampled before assignment.  The exact construction-side bounds are
\[
 \inf_{\widehat\theta}
 \mathbb E_{\overline Q_b^{\mathrm{blk}},\pi}
 [ (\widehat\theta-\theta_T(Y))^2]
 \ge \frac{B^2}{576}
 \mathbb E_{\overline Q_b^{\mathrm{blk}},\pi}
 [\overline{\mathsf V}_b(W,R_b)]
\]
and the common-null floor is \(B^2/(12T)\).  The band calculations give
\[
D_{\mathrm{KL}}\!\left(
P_{Q_0^{T,L},\pi}\|P_{Q_{a,b}^{\mathrm{band}},\pi}
\right)
\le \frac{16A_L^2}{3B^2}
\mathbb E_{Q_0^{T,L},\pi}[\mathsf I_{a,b}(W)],
\qquad
\Delta\mathbb E[\theta_T]=\frac{2A_L}{3}.
\]
In particular, these imply the explicit constrained-risk lower bound
\[
 \frac{A_L^2}{72}
 \left[
 1-\left\{
 \frac{8A_L^2}{3B^2}
 \mathbb E_{Q_0^{T,L},\pi}[\mathsf I_{a,b}(W)]
 \right\}^{1/2}
 -\frac{24B^2}{TA_L^2}
 \right]_+
\]
for the maximum risk over \(Q_0^{T,L}\) and \(Q_{a,b}^{\mathrm{band}}\).
The assignment-kernel factors cancel exactly in both structural likelihood ratios.  These equations rigorously settle the support, posterior-variance, separation, and band-information parts of the proposed converse.  Only the policy-uniform dyadic charging and common-law transfer remain open; consequently neither the logarithmic certificate nor removal of its logarithm is proved.

%%% FIELD frontier-proof
This is a conditional implication, and we use exactly the no-loss premise stated in the theorem.  Put \(D=B\wedge C\), \(p=(2\beta+1)^{-1}\), and \(x=T^p\).  The assumed threshold gives \(L_T\le T/4\).  For a fixed \(\pi\in\mathcal D_T\), the assumed no-loss certificate selects, before assignment, a prior \(Q_{\pi,L_T}\) supported on \(\mathcal P_{T,\beta}(B,C)\) such that
\[
 \inf_{\widehat\theta}
 \mathbb E_{Q_{\pi,L_T},\pi}
 [ (\widehat\theta-\theta_T(Y))^2]
 \ge \kappa_\beta
 \min\left\{\frac{B^2L_T}{T},\ B^2,\ C^2L_T^{-2\beta}\right\}.
\]
For every estimator, its supremum risk over the class is at least its average risk under a prior supported on the class.  Therefore
\[
 \inf_{\widehat\theta}\sup_{Y\in\mathcal P_{T,\beta}(B,C)}
 \mathbb E_{\pi,Y}[(\widehat\theta-\theta_T(Y))^2]
 \ge \kappa_\beta
 \min\left\{\frac{B^2L_T}{T},\ B^2,\ C^2L_T^{-2\beta}\right\}.
\]
The elementary floor bounds \(x/2\le L_T\le x\), valid because \(x\ge1\), imply
\[
 \frac{L_T}{T}\ge\frac12T^{-2\beta/(2\beta+1)},
 \qquad
 L_T^{-2\beta}\ge T^{-2\beta/(2\beta+1)}.
\]
Also \(T^{-2\beta/(2\beta+1)}\le1\).  Hence
\[
 \min\left\{\frac{B^2L_T}{T},\ B^2,\ C^2L_T^{-2\beta}\right\}
 \ge \frac12\min\{B^2,C^2\}
 T^{-2\beta/(2\beta+1)}.
\]
Because \(\min\{B^2,C^2\}=\min\{B^2,(B\wedge C)^2\}\), and because the bound holds for every \(\pi\), taking both infima in \(\mathrm{def{:}minimax\mbox{-}mse}\) proves the asserted lower inequality.  The upper inequality is the radius-free corollary of \(\mathrm{thm{:}radius\mbox{-}aware\mbox{-}oracle\mbox{-}upper}\).  Positivity of \(B\), \(C\), and \(\kappa_\beta\), and finiteness of \(K_+(B,C,\beta)\), give the conditional \(\Theta_{B,C,\beta}\) conclusion.  No step of this proof establishes the no-loss certificate; that premise remains open in \(\mathrm{oeq{:}sequential\mbox{-}multiscale\mbox{-}certificate}\).

%%% FIELD bojinov-scope-statement
Bojinov, Simchi-Levi, and Zhao study a known exact finite carryover order, optimize Horvitz--Thompson risk over a precommitted class of regular switchback designs, and use a treatment-effect estimand with an initial boundary truncation.  Their result is not an all-sequential-policy minimax theorem over arbitrary estimators.

%%% FIELD upper-proposed-statement
For every \(T\ge2\), \(C\ge0\), and known \(\beta>0\), the radius-aware certified window \(b_\star\) satisfies
\[
 R_T^{(2)}(\beta;B,C)
 \le
 \min\!\left\{
 4B^2,
 \min_{1\le b\le T}
 \left[
 B\sqrt{\frac{(8\mathrm e-4)b}{T}}
 +\frac{2\mathrm e C}{\beta}(b+1)^{-\beta}
 \right]^2
 \right\}.
\]
Writing
\[
 x=\left\{
 \frac{(2\mathrm e C/\beta)^2T}{(8\mathrm e-4)B^2}
 \right\}^{1/(2\beta+1)},
 \qquad z=1\vee(x\wedge T),
\]
one has the explicit scale comparison
\[
 \frac12(1+\sqrt2)^{-1/\beta}z
 \le b_\star\le(1+\sqrt2)^2z.
\]
Consequently,
\[
 b_\star\asymp_\beta
 1\vee\left(
 \left\{(C/B)^2T\right\}^{1/(2\beta+1)}\wedge T
 \right).
\]
If \(x\le1\), then
\[
 R_T^{(2)}(\beta;B,C)
 \le
 \min\left\{4B^2,(8\mathrm e-4)(1+2^{-\beta})^2\frac{B^2}{T}\right\}.
\]
If \(1\le x\le T\), then
\[
 R_T^{(2)}(\beta;B,C)
 \le
 \min\left\{4B^2,
 K_{\mathrm{mid},\beta}
 B^{4\beta/(2\beta+1)}C^{2/(2\beta+1)}
 T^{-2\beta/(2\beta+1)}
 \right\},
\]
where
\[
 K_{\mathrm{mid},\beta}
 =(1+\sqrt2)^2(8\mathrm e-4)^{2\beta/(2\beta+1)}
 (2\mathrm e/\beta)^{2/(2\beta+1)}.
\]
If \(x\ge T\), then
\[
 R_T^{(2)}(\beta;B,C)
 \le
 \min\left\{4B^2,
 \left[B\sqrt{8\mathrm e-4}
 +\frac{2\mathrm e C}{\beta}(T+1)^{-\beta}\right]^2
 \right\}.
\]
Finally,
\[
 R_T^{(2)}(\beta;B,C)
 \le K_+(B,C,\beta)T^{-2\beta/(2\beta+1)}.
\]
This theorem is a finite-horizon polynomial-envelope upper certificate for the all-time GATE, delivered by the stated Markov design and truncated Horvitz--Thompson estimator.  For comparison, Bojinov, Simchi-Levi, and Zhao optimize Horvitz--Thompson risk for known exact finite carryover over precommitted regular designs and a different boundary-truncated estimand.  Neither that comparison nor the present upper certificate asserts all-policy optimality or an arbitrary-estimator converse.

%%% FIELD upper-proof
Put \(A=\sqrt{8\mathrm e-4}\), \(D=2\mathrm e/\beta\), \(a=AB/\sqrt T\), \(c=DC\), and
\[
 f(b)=a\sqrt b+c(b+1)^{-\beta},
 \qquad b\in\{1,\ldots,T\}.
\]
Here \(a>0\) and \(c\ge0\).  By \(\mathrm{lem{:}markov\mbox{-}upper\mbox{-}audit}\), the design \(\pi^{(b)}\) and estimator \(\widehat\theta_b\) have worst-case mean squared error at most \(f(b)^2\).  Thus \(\mathrm{def{:}minimax\mbox{-}mse}\) and \(\mathrm{def{:}oracle\mbox{-}scale}\) give
\[
 R_T^{(2)}(\beta;B,C)\le\min_{1\le b\le T}f(b)^2=f(b_\star)^2.
\]
Also, boundedness and \(\mathrm{def{:}gate}\) imply
\[
 |\theta_T(Y)|
 \le T^{-1}\sum_{t=1}^T
 \{|Y_t(\mathbf1_{1:t})|+|Y_t(\mathbf0_{1:t})|\}
 \le2B.
\]
The identically zero estimator therefore has squared error at most \(4B^2\) for every array and every design.  Combining the two available designs and estimators proves the first display.

We prove the scale assertion.  If \(c>0\), the displayed \(x\) satisfies
\[
 x=(c/a)^{1/(\beta+1/2)},
 \qquad c=ax^{\beta+1/2};
\]
if \(c=0\), set \(x=0\), as in the displayed formula.  Put \(M=1+\sqrt2\).  If \(x\le1\), optimality and evaluation at \(b=1\) give
\[
 a\sqrt{b_\star}\le f(b_\star)\le f(1)
 =a+c2^{-\beta}\le a(1+2^{-\beta}).
\]
Thus \(1\le b_\star\le(1+2^{-\beta})^2\le M^2\).  Since \(z=1\), both asserted scale inequalities follow.

Suppose \(1<x<T\), and let \(m=\min\{T,\lceil x\rceil\}\).  Then \(m\le2x\) and \(m+1\ge x\), whence
\[
 f(m)\le a\sqrt{2x}+cx^{-\beta}=Ma\sqrt x.
\]
The nonnegative variance term and optimality imply \(b_\star\le M^2x\).  The nonnegative bias term similarly gives
\[
 c(b_\star+1)^{-\beta}
 \le Ma\sqrt x,
 \qquad
 b_\star+1\ge M^{-1/\beta}x.
\]
If \(M^{-1/\beta}x\le2\), then \(b_\star\ge1\ge M^{-1/\beta}x/2\); otherwise \(b_\star\ge M^{-1/\beta}x-1\ge M^{-1/\beta}x/2\).  This proves the intermediate comparison.

If \(x\ge T\), evaluation at \(b=T\) yields
\[
 c(b_\star+1)^{-\beta}
 \le a\sqrt T+cT^{-\beta}.
\]
Writing \(u=x/T\ge1\) and using \(c=ax^{\beta+1/2}\), we obtain
\[
 (b_\star+1)^\beta
 \ge T^\beta\frac{u^{\beta+1/2}}{1+u^{\beta+1/2}}
 \ge T^\beta/2.
\]
Since \(M>2\), the same two-case argument gives \(b_\star\ge M^{-1/\beta}T/2\), while \(b_\star\le T\le M^2T\) is automatic.  The three regimes prove the explicit comparison, and replacing \(x\) by its formula gives the \(\asymp_\beta\) statement.

For the low-memory risk bound, choose \(b=1\).  The identity \(c=ax^{\beta+1/2}\) and \(x\le1\) yield
\[
 f(1)^2\le a^2(1+2^{-\beta})^2
 =(8\mathrm e-4)(1+2^{-\beta})^2B^2/T.
\]
For \(1\le x\le T\), the preceding integer \(m\) satisfies \(f(m)^2\le M^2a^2x\), and direct substitution gives
\[
 M^2a^2x
 =K_{\mathrm{mid},\beta}
 B^{4\beta/(2\beta+1)}C^{2/(2\beta+1)}
 T^{-2\beta/(2\beta+1)}.
\]
For \(x\ge T\), evaluation at \(b=T\) gives exactly the second entry in the stated minimum.  The zero-estimator argument supplies the first entry in every regime.

For the radius-free corollary, put \(p=(2\beta+1)^{-1}\) and \(b_0=\lceil T^p\rceil\).  Since \(T\ge2\),
\[
 b_0\le T^p+1\le T^p(1+2^{-p}),
 \qquad
 (b_0+1)^{-\beta}\le T^{-\beta p}.
\]
The optimized certificate evaluated at \(b_0\) gives
\[
 \begin{aligned}
 R_T^{(2)}(\beta;B,C)
 &\le\left\{
 B\sqrt{(8\mathrm e-4)(1+2^{-p})}
 +\frac{2\mathrm e C}{\beta}
 \right\}^2T^{-2\beta p}\\
 &=K_+(B,C,\beta)T^{-2\beta/(2\beta+1)}.
 \end{aligned}
\]
The theorem's final scope comparison is the cited content of \(\mathrm{lem{:}bojinov\mbox{-}regular\mbox{-}design\mbox{-}scope}\); none of the inequalities above uses that comparator as a converse.
